\section{Eventual coverage of component counts}
\label{sec:scalar-finality}

This section is the scalar shadow of the categorical transfer.  It retains only the width of the component-count interval and works fibrewise over the full prepared-start and sponsor parameter objects.

Let $0<\sigma<1$ and $0<\rho<\kappa\sigma\rho_P$ be arbitrary parameters as in Section~\ref{sec:local-profile}.  Theorem~\ref{thm:uniform-profile} leaves two scalar quantities at a prime-power step: the required incoming width $\delta_\rho(Q)$ and, at prime steps, the positive width increment $\mu_\rho(q)$.

Enumerate the prime powers
\[
\{Q=p^e:Q\ge Q_{\mathrm{tr}}\}
\]
in increasing order.  For an integer interval $I=[A,B]$ write $w(I)=B-A$.  At a prime $q$, Theorem~\ref{thm:uniform-profile} gives
\begin{equation}\label{eq:prime-action}
[A,B]\longmapsto
[A+\delta_\rho(q),\,B+\delta_\rho(q)+\mu_\rho(q)],
\end{equation}
so the width increases by $\mu_\rho(q)$.  At a composite prime power $Q=p^e$, $e\ge2$, it gives
\begin{equation}\label{eq:composite-action}
[A,B]\longmapsto[A+\delta_\rho(Q),B],
\end{equation}
so the width decreases by $\delta_\rho(Q)$.  A step is available whenever the incoming width is at least $\delta_\rho(Q)$; then the old and new integer intervals overlap.

Put
\[
\Phi(X)=\frac X{\log X}.
\]
For every $X\ge1$ define
\[
\Delta(X)=
\max\Bigl(\{0\}\cup
\{\delta_\rho(Q):Q_{\mathrm{tr}}\le Q\le X,\ Q\text{ a prime power}\}\Bigr),
\]
and
\begin{equation}\label{eq:cumulative-debt}
\mathcal L(X)=
\sum_{\substack{Q_{\mathrm{tr}}\le p^e\le X\\ e\ge2}}\delta_\rho(p^e).
\end{equation}
Thus both functions are defined on all positive integers and vanish below the transfer threshold.  By \eqref{eq:d-composite},
\begin{equation}\label{eq:max-demand-subscale}
\Delta(X)=O_\rho((\log X)^2)=o(\Phi(X)).
\end{equation}
The number of prime powers $p^e\le X$ with $e\ge2$ is $O(\sqrt X\log X)$, so
\begin{equation}\label{eq:debt-subscale}
\mathcal L(X)=O_\rho(\sqrt X\,\log^3X)=o(\Phi(X)).
\end{equation}

\subsection{Bounded-lag prime compensation as a sponsor object}

Let $\Lambda\ge2$ be the multiplier in Proposition~\ref{prop:prime-supply-main} and put
\begin{equation}\label{eq:sponsor-lag}
\Lambda_s=4\Lambda.
\end{equation}
For a sufficiently large prime-power rank $Q$ define the finite sponsor object
\begin{equation}\label{eq:sponsor-bounded-lag}
\operatorname{Sp}(Q)=
\{q\ge Q_{\mathrm{tr}}:q\text{ prime},\ q<Q\le\Lambda_s q\}.
\end{equation}
Proposition~\ref{prop:prime-supply-main} implies that $\operatorname{Sp}(Q)$ is nonempty for every sufficiently large $Q$: with
\[
Y=\left\lfloor\frac{Q}{2\Lambda}\right\rfloor,
\]
the prime band $(Y,\Lambda Y]$ maps into \eqref{eq:sponsor-bounded-lag} after a fixed onset.

Define the lower sponsor gain by the universal lower envelope
\[
\underline\mu_\rho(Q)=
\min_{q\in\operatorname{Sp}(Q)}\mu_\rho(q).
\]
Every $q\in\operatorname{Sp}(Q)$ satisfies $Q/\Lambda_s\le q<Q$.  Hence \eqref{eq:mint-prime} gives
\begin{equation}\label{eq:sponsor-mint}
\underline\mu_\rho(Q)
\gg_{\rho,\Lambda_s}\frac Q{\log Q}
=\Theta_{\rho,\Lambda_s}(\Phi(Q)).
\end{equation}
The compensation estimate is therefore a property of the whole sponsor object.

\subsection{Prepared starts and continuation}

Define
\begin{equation}\label{eq:startup-request}
R_\rho(B)=\Delta(\Lambda_sB)+\mathcal L(\Lambda_sB).
\end{equation}
By \eqref{eq:max-demand-subscale} and \eqref{eq:debt-subscale},
\begin{equation}\label{eq:startup-subscale-scalar}
R_\rho(B)=o(B/\log B),
\end{equation}
so $R_\rho\in\mathfrak L$.

Let $0<\eta<\eta_0$ be arbitrary.  Proposition~\ref{prop:wide-start} says that the set of $B$ for which
\[
\mathfrak W_B(\eta,\sigma,R_\rho)\ne\varnothing
\]
contains a final segment.  For every $\xi\in\mathfrak W_B(\eta,\sigma,R_\rho)$ write
\[
R_B^\xi:\{0_M\}\rightsquigarrow
\operatorname{Fib}_{E_B}(\gamma_B^\xi)\times I_B^\xi
\]
for its covering correspondence; the interval $I_B^\xi$ has width $R_\rho(B)$.

Starting from $I_B^\xi$, apply the scalar interval maps \eqref{eq:prime-action} and \eqref{eq:composite-action} in increasing prime-power order.  Let $I_X^\xi$ denote the interval after all steps of rank at most $X$ have been processed whenever those steps are available.

Define $\mathcal B_{\mathrm{cont}}(\rho)$ to be the set of integers $B$ for which, for every admissible $\eta,\sigma$, every $\xi\in\mathfrak W_B(\eta,\sigma,R_\rho)$ and every prime-power step $Q>B$, all steps through $Q$ are available and, whenever $Q>\Lambda_sB$,
\begin{equation}\label{eq:sponsored-width}
w(Q)\ge\underline\mu_\rho(Q)-\mathcal L(Q).
\end{equation}

\begin{lemma}[Continuation]\label{lem:component-count-continuation}
The set $\mathcal B_{\mathrm{cont}}(\rho)$ contains a final segment of $\N$.
\end{lemma}

\begin{proof}
The proof is independent of the prepared start because only the initial width $R_\rho(B)$ is used.  For $B<Q\le\Lambda_sB$, discarding all positive prime increments gives
\[
w(Q)\ge
R_\rho(B)-\mathcal L(Q)
\ge
\Delta(\Lambda_sB)
\ge
\delta_\rho(Q),
\]
so every step in the initial transition is available.

For $Q>\Lambda_sB$, all $q\in\operatorname{Sp}(Q)$ satisfy $q>Q/\Lambda_s>B$.  Once the preceding steps are available, all sponsor prime steps have therefore occurred.  The current width contains every accumulated positive prime gain and loses at most $\mathcal L(Q)$ through composite prime powers, so
\[
w(Q)\ge\mu_\rho(q)-\mathcal L(Q)
\qquad(q\in\operatorname{Sp}(Q)).
\]
Taking the minimum over the sponsor object gives \eqref{eq:sponsored-width}.  By \eqref{eq:sponsor-mint} and \eqref{eq:debt-subscale}, its right-hand side is $\gg\Phi(Q)$ for large $Q$, whereas \eqref{eq:max-demand-subscale} gives $\delta_\rho(Q)=o(\Phi(Q))$.  Hence all sufficiently large $B$ belong to $\mathcal B_{\mathrm{cont}}(\rho)$.
\end{proof}

For $B\in\mathcal B_{\mathrm{cont}}(\rho)$ and every prepared start $\xi$, successive intervals overlap by construction and their upper endpoints are unbounded because the positive prime increments tend to infinity.  Thus the increasing union of the intervals $I_X^\xi$ contains a final ray in $\N$.

\subsection{Vanishing of the affine centre and cutoff incidence}

For $X\ge B$ let
\[
\pi_{B,X}:A/E_B\longrightarrow A/E_X
\]
be the quotient map and put
\[
\gamma_X^\xi=\pi_{B,X}(\gamma_B^\xi).
\]
The quotient maps form the commuting system
\begin{equation}\label{eq:centre-transport-square}
\begin{CD}
A/E_B @>{\pi_{B,X}}>> A/E_X\\
@V{\pi_{B,Y}}VV @VV{\pi_{X,Y}}V\\
A/E_Y @= A/E_Y
\end{CD}
\qquad(B\le X\le Y).
\end{equation}
By \eqref{eq:centre-kernel-union}, the vanishing set
\[
\mathcal V_\xi=
\{X\ge B:\gamma_X^\xi=0\}
\]
is a nonempty final segment of $\N$.

Let
\[
\mathcal H_B=\{X\in\N:X\ge B\}
\]
and define the cutoff--count incidence object
\begin{equation}\label{eq:term-k}
\mathfrak X_\xi=
\{(k,X)\in\N_{>0}\times\mathcal H_B:
X\in\mathcal V_\xi,\ k\in I_X^\xi\},
\end{equation}
with projection
\[
p_\xi:\mathfrak X_\xi\longrightarrow\N_{>0},
\qquad
(k,X)\longmapsto k.
\]
Retain the whole object of admissible scalar thresholds
\begin{equation}\label{eq:scalar-threshold-object}
\mathfrak K_\xi=
\{k_0\in\N_{>0}:[k_0,\infty)\cap\N\subseteq\operatorname{im}(p_\xi)\}.
\end{equation}

Define the scalar-good start object
\begin{equation}\label{eq:scalar-good-start-object}
\mathfrak G_{\mathrm{sc}}(\eta,\sigma,\rho)
=
\{(B,\xi):B\in\mathcal B_{\mathrm{cont}}(\rho),\
\xi\in\mathfrak W_B(\eta,\sigma,R_\rho),\
\mathfrak K_\xi\ne\varnothing\},
\end{equation}
with projection $b_{\mathrm{sc}}(B,\xi)=B$.

\begin{theorem}[Scalar finality]\label{thm:scalar-terminality}
For every admissible $(\eta,\sigma,\rho)$, the regular image of
\[
b_{\mathrm{sc}}:
\mathfrak G_{\mathrm{sc}}(\eta,\sigma,\rho)
\longrightarrow\N
\]
is cofinite.  Equivalently, every sufficiently large $B$ supports prepared starts whose cutoff--count projection has cofinite image; indeed every prepared start over such a $B$ has this property.
\end{theorem}

\begin{proof}
Proposition~\ref{prop:wide-start} and Lemma~\ref{lem:component-count-continuation} show that, for all sufficiently large $B$, the prepared-start object is nonempty and $B\in\mathcal B_{\mathrm{cont}}(\rho)$.  For every $\xi$ over such a $B$, the vanishing set $\mathcal V_\xi$ is a final segment, while the successive-overlap sequence $(I_X^\xi)$ has unbounded upper endpoints.  Hence
\[
\operatorname{im}(p_\xi)
=
\bigcup_{X\in\mathcal V_\xi}I_X^\xi
\]
contains a final ray, which is exactly the nonemptiness of \eqref{eq:scalar-threshold-object}.  Thus every sufficiently large $B$ lies in $\operatorname{im}(b_{\mathrm{sc}})$.
\end{proof}
